% Base mala frente a base buena, con la celosia del Ejemplo 1 del enunciado.
%   (a) Base publica (mala):   (5, 19), (3, 14)
%   (b) Base reducida (buena): (3, 1), (-1, 4)   (LLL la devuelve con el signo cambiado)
% Las dos generan la MISMA celosia: los puntos son iguales en los dos dibujos.
% Lo unico que cambia es el paralelogramo fundamental (sombreado): largo y estrecho
% con la base mala, casi un cuadrado con la buena.
\begin{tikzpicture}[x=0.155cm,y=0.155cm,font=\footnotesize,
                    vec/.style={-{Stealth[length=3.5pt]},line width=0.8pt}]
  \newcommand{\puntos}{%
    \foreach \a in {-8,...,10}{\foreach \b in {-3,...,6}{
      \fill ({3*\a-\b},{\a+4*\b}) circle (0.38);}}}

  % (a) Base publica
  \begin{scope}
    \clip (-3,-2) rectangle (12,22);
    \fill[gray!25] (0,0) -- (5,19) -- (8,33) -- (3,14) -- cycle;
    \puntos
    \draw[vec] (0,0) -- (5,19);
    \draw[vec] (0,0) -- (3,14);
  \end{scope}
  \draw[line width=0.4pt] (-3,-2) rectangle (12,22);
  \node[anchor=west] at (5.3,19.5) {$b_1$};
  \node[anchor=east] at (2.7,14.5) {$b_2$};
  \node[below] at (4.5,-2) {(a) Base pública};

  % (b) Base reducida
  \begin{scope}[shift={(19,0)}]
    \begin{scope}
      \clip (-3,-2) rectangle (12,22);
      \fill[gray!25] (0,0) -- (3,1) -- (2,5) -- (-1,4) -- cycle;
      \puntos
      \draw[vec] (0,0) -- (3,1);
      \draw[vec] (0,0) -- (-1,4);
    \end{scope}
    \draw[line width=0.4pt] (-3,-2) rectangle (12,22);
    \node[anchor=west] at (3.2,1) {$b_1'$};
    \node[anchor=east] at (-1.1,4.3) {$b_2'$};
    \node[below] at (4.5,-2) {(b) Base reducida};
  \end{scope}
\end{tikzpicture}
